% TOP_PROBLEM: {"problemId": "problem.calderon-uniqueness-for-bounded-measurable-isotropic-conductivities-in-higher-dimensions", "canonicalTitle": "Calderón uniqueness for bounded measurable real isotropic conductivities", "releaseRank": 314, "primaryDomain": "analysis_pde", "primaryDomainLabel": "Analysis & PDE", "displayStatus": "open_with_solved_subcases", "releaseStatus": "open_with_solved_subcases"}
% !TeX program = lualatex
% ATTEMPT_STATUS: unresolved
% Model: GPT-6 Astra Ultra; continuation dated 27 September 2026.
\documentclass[11pt]{article}
\usepackage[margin=25mm]{geometry}
\usepackage{fontspec}
\setmainfont{Latin Modern Roman}
\usepackage{amsmath,amssymb,mathtools}
\usepackage{unicode-math}
\setmathfont{Latin Modern Math}
\usepackage{xurl,hyperref}
\hypersetup{colorlinks=true,urlcolor=blue}
\setcounter{secnumdepth}{0}
\setlength{\parindent}{0pt}
\setlength{\parskip}{5pt}
\setlength{\emergencystretch}{3em}
\begin{document}
\section*{\texorpdfstring{Calderón uniqueness for bounded measurable real isotropic conductivities}{Calderón uniqueness for bounded measurable real isotropic conductivities}}\label{314}
\subsection{Definitions and mathematical statement}
Let $n\ge3$ and let $\Omega\subset{\mathbb{R}}^n$ be a bounded connected Lipschitz domain. For a real scalar conductivity $\gamma\in L^\infty(\Omega)$ with $0<c\le\gamma\le C<\infty$, the weak solution with boundary value $f\in H^{1/2}(\partial\Omega)$ satisfies $\nabla\cdot(\gamma\nabla u)=0$. Its Dirichlet-to-Neumann map is the bounded operator determined by
\[
 \langle\Lambda_\gamma f,g\rangle
 =\int_\Omega\gamma\nabla u\cdot\nabla v{\,\mathrm{d}} x,
 \qquad v|_{\partial\Omega}=g.
\]
The value is independent of the chosen $H^1$ extension $v$. The uniqueness question is whether $\Lambda_{\gamma_1}=\Lambda_{\gamma_2}$ implies $\gamma_1=\gamma_2$ almost everywhere. The hypotheses impose no derivative regularity, smallness, or prior agreement near the boundary. This scalar question is distinct from anisotropic uniqueness modulo boundary-fixing changes of coordinates.
\par\smallskip\textit{Formulation and concept references:} \href{https://arxiv.org/pdf/1803.06931v2}{[S1]}.

\subsection{Short English statement}
Do all boundary voltage-to-current measurements uniquely determine a bounded measurable positive scalar conductivity in dimension at least three?
\subsection{Research attempt}
\textbf{Endpoint and source audit.}
The coefficients here are real scalar $L^\infty$ functions with a fixed
positive lower bound.  There is no Sobolev derivative available and no
assumption that two coefficients agree in a collar of the boundary.
The primary note [S1] proposes a Clifford-algebra/Beltrami strategy as a
possible first step toward this endpoint; its abstract does not claim the
full measurable-conductivity uniqueness theorem.  Results for Lipschitz or
other derivative-controlled conductivities and two-dimensional measurable
results do not remove these distinctions.  We give exact identities, two
rough-coefficient subcases, and injectivity of the linearization at a constant
coefficient.  The final nonlinear completeness step remains unproved.

\textbf{The variational formulation needs no coefficient derivatives.}
Fix a bounded extension operator from $H^{1/2}(\partial\Omega)$ to
$H^1(\Omega)$, available for a Lipschitz domain.  If $F$ extends $f$, seek
$u=F+w$, $w\in H_0^1(\Omega)$, with
\[
 \int_\Omega\gamma\nabla w\cdot\nabla v
 =-\int_\Omega\gamma\nabla F\cdot\nabla v
 \quad(v\in H_0^1(\Omega)).
\]
Boundedness, uniform ellipticity, and Poincar\'e's inequality give existence
and uniqueness by Lax--Milgram.  Testing with $w$ gives a bound for
$\|u\|_{H^1}$ in terms of $f$, the domain, and the ellipticity constants.
Changing an extension of the second boundary argument by an element of
$H_0^1$ changes the DN pairing by zero.  This proves the asserted
well-definedness and boundedness of $\Lambda_\gamma$ without assigning a
classical normal derivative to $u$.

Let $u_1^f$ solve for $\gamma_1$ with trace $f$, and $u_2^g$ solve for
$\gamma_2$ with trace $g$.  Symmetry of each real-coefficient DN form and
testing the weak equations with differences having zero trace give
\[
 \langle(\Lambda_{\gamma_1}-\Lambda_{\gamma_2})f,g\rangle
 =\int_\Omega(\gamma_1-\gamma_2)
                    \nabla u_1^f\cdot\nabla u_2^g\,dx.
\]
For example the first term can use $u_2^g$ as its extension; the second can
use $u_1^f$ after symmetry and the zero-trace replacement of its own solution.
This is the basic orthogonality identity.  Equality of the DN maps supplies
orthogonality to these products, not orthogonality to every $L^1$ function.
The density or special-solution theorem needed to bridge that gap is the
substantive issue.

\textbf{Monotonicity and an ordered constant-background subcase.}
For a fixed real boundary value $f$, the solution uniquely minimizes
\[
 E_\gamma(v)=\int_\Omega\gamma|\nabla v|^2
\]
among all $v$ with trace $f$.  Indeed writing $v=u+w$ makes the cross term
zero and leaves $E_\gamma(v)=E_\gamma(u)+E_\gamma(w)$.
For solutions $u_1,u_2$ with that common trace, comparison with the other
minimizer gives
\[
 \int(\gamma_1-\gamma_2)|\nabla u_1|^2
 \leq\langle(\Lambda_{\gamma_1}-\Lambda_{\gamma_2})f,f\rangle
 \leq\int(\gamma_1-\gamma_2)|\nabla u_2|^2.
\]
In particular pointwise ordering of conductivities implies ordering of their
DN quadratic forms.

Suppose now that $\gamma\geq c_0>0$ a.e., where $c_0$ is constant, and
$\Lambda_\gamma=\Lambda_{c_0}$.  Use $f=x_j|_{\partial\Omega}$.
The constant-coefficient solution is $x_j$.  If $u$ is the $\gamma$ solution,
\[
 E_{c_0}(x_j)\leq E_{c_0}(u)\leq E_\gamma(u)
 =\langle\Lambda_\gamma f,f\rangle=E_{c_0}(x_j).
\]
Equality and uniqueness of the constant-coefficient minimizer force $u=x_j$.
It follows that $\int_\Omega(\gamma-c_0)=0$, hence $\gamma=c_0$ a.e.
If instead $\gamma\leq c_0$, use
\[
 E_\gamma(u)\leq E_\gamma(x_j)=\int\gamma
 \leq c_0|\Omega|=E_{c_0}(x_j)=E_\gamma(u)
\]
to reach the same conclusion.  This proves uniqueness against a constant
background under a one-sided ordering, with no regularity assumption on
$\gamma$.  It does not prove uniqueness for two unordered coefficients.
For a general variable comparison coefficient, the corresponding gradients
are unknown; declaring them nonzero on every possible discrepancy set would
insert an additional unique-continuation or localization assertion.

\textbf{A measurable layered subcase with full boundary data.}
Take $\Omega=(0,1)^n$ and $\gamma(x)=a(x_1)$, where $a$ is merely measurable
and $0<c\leq a\leq C$.  For the boundary voltage $f=x_2$, the function
$u=x_2$ is a weak solution because
$\operatorname{div}(a(x_1)e_2)=0$ in distributions.
For a smooth boundary test $g$ compactly supported in the relative interior
of the face $\{x_2=0\}$, the variational DN pairing is
\[
 \langle\Lambda_a f,g\rangle
 =-\int_{(0,1)^{n-1}}a(x_1)g(x_1,0,x_3,\ldots,x_n)
                                      \,dx_1dx_3\cdots dx_n.
\]
To justify the formula at this regularity, take a smooth extension supported
near that face and integrate $a(x_1)\partial_2v$ in the $x_2$ variable using
Fubini.  No derivative or boundary trace of $a$ in its own variable is taken.

If $a_1,a_2$ have equal full DN maps, choose product tests
$g=\psi(x_1)\chi(x_3,\ldots,x_n)$ with $\chi$ of integral one.
Then $\int_0^1(a_1-a_2)\psi=0$ for every compactly supported smooth $\psi$,
so $a_1=a_2$ a.e.  This subcase genuinely permits jumps and arbitrarily
oscillatory measurable layers.  It uses accessible side faces and the a priori
layered geometry, not just measurements on the two faces perpendicular to
the layers.  Those more limited measurements can collapse to an effective
conductivity and would not justify this argument.

\textbf{Injectivity of the derivative at the identity conductivity.}
Let $\gamma_\varepsilon=1+\varepsilon h$ with real
$h\in L^\infty(\Omega)$ and $|\varepsilon|\|h\|_\infty\leq1/2$.
Write $u_\varepsilon^f=u_0^f+w_\varepsilon$, where $u_0^f$ is harmonic.
The weak equation gives
\[
 \int(1+\varepsilon h)\nabla w_\varepsilon\cdot\nabla v
 =-\varepsilon\int h\nabla u_0^f\cdot\nabla v,
 \quad
 \|\nabla w_\varepsilon\|_2
 \leq2|\varepsilon|\|h\|_\infty\|\nabla u_0^f\|_2.
\]
The preceding integral identity therefore yields the Fr\'echet derivative
\[
 \langle D\Lambda_1[h]f,g\rangle
 =\int_\Omega h\nabla u_0^f\cdot\nabla u_0^g,
\]
with a remainder bounded by a constant times
$\varepsilon^2\|h\|_\infty^2\|f\|_{H^{1/2}}\|g\|_{H^{1/2}}$.
Complexify this real bilinear identity.  For $\xi\ne0$, choose a real unit
vector $e\perp\xi$ and put
\[
 \zeta=\frac{i\xi}{2}+\frac{|\xi|e}{2},\qquad
 \eta=\frac{i\xi}{2}-\frac{|\xi|e}{2}.
\]
Then $\zeta\cdot\zeta=\eta\cdot\eta=0$,
$\zeta+\eta=i\xi$, and $\zeta\cdot\eta=-|\xi|^2/2$.
The functions $e^{\zeta\cdot x}$ and $e^{\eta\cdot x}$ are smooth harmonic
functions on the bounded domain.  If $D\Lambda_1[h]=0$, substitution gives
\[
 0=-\frac{|\xi|^2}{2}\int_\Omega h(x)e^{i\xi\cdot x}\,dx.
\]
The zero extension of $h$ is integrable and compactly supported.  Its Fourier
transform vanishes at every nonzero frequency and, by continuity, at zero.
Fourier uniqueness gives $h=0$ a.e.  Thus the linearized inverse problem is
injective throughout the requested $L^\infty$ class of perturbation directions.

\textbf{The nonlinear gap is not removed by this derivative.}
An injective derivative in infinite-dimensional spaces need not have a
bounded inverse in the relevant norms and does not itself imply nonlinear
global injectivity.  At a variable coefficient, the exponentials used above
are not conductivity solutions.  Replacing them by complex geometric optics
solutions with errors small enough to recover Fourier data is a new analytic
construction, not a formal repetition of the constant-coefficient proof.
Likewise the formal substitution $v=\sqrt\gamma\,u$ introduces
$q=\gamma^{-1/2}\Delta\sqrt\gamma$.  At mere $L^\infty$ regularity this
expression cannot automatically be treated as an ordinary potential or
multiplied with arbitrary rough solutions.  Smooth approximation may produce
unbounded derivative norms: $1+\delta\sin(jx_1)$ stays uniformly elliptic
for fixed $0<\delta<1$ while its gradient grows with $j$.

Finally, a boundary-fixing change of variables transforms a scalar coefficient
into a tensor proportional to $DF\,DF^T/|\det DF|$, composed with the inverse
map.  That tensor is generally not scalar.  The usual anisotropic gauge
therefore does not directly provide a counterexample here, and singular
cloaking coefficients violate the stipulated bounded positive scalar class.

\textbf{Verification and outcome.}
The diagnostic checks the complex-vector identities numerically at specified
frequencies and checks layered boundary pairings exactly for a piecewise
constant rational example.  The infinite-dimensional variational and Fourier
arguments are supplied above, not inferred from that finite sample.  The
attempt proves ordered constant-background uniqueness, layered uniqueness,
and injectivity of the constant-background linearization.  It supplies no
complete family of products of solutions for two arbitrary measurable
coefficients, which is the first missing step toward the full target.
\subsection{Sources}
\begin{itemize}
\item[S1]Matteo Santacesaria, Note on Calderón's inverse problem for measurable conductivities, Inverse Problems and Imaging13(1)(2019),149–157; arXiv1803.06931v2.
\url{https://arxiv.org/pdf/1803.06931v2}.
\textit{Formulation source.}
\end{itemize}
\end{document}
